
% \tdplotsetmaincoords{80}{115}
\begin{tikzpicture}[node distance=0.8cm and 0.5cm, font=\sffamily , perspective={
    p = {(0,0,0)}}]

% \begin{tikzpicture}[node distance=0.8cm and 0.5cm, font=\sffamily, tdplot_main_coords ]

  % Styles for boxes, ellipses, and colored labels



  % Gene nodes and numbers for container type 1
  \node[gene_node ] (order1) {Order};
  \node[label_node, fill=orange, above=0.25cm of order1] (o1) {0.5};

  \node[gene_node , right=\separationofcircles cm of order1] (orientation1) {Orientation};
  \node[label_node=lime, above=0.25cm of orientation1]  {0.3};

  \node[gene_node , right=\separationofcircles cm of orientation1] (layersXY1) {Layers \\ xy plane};
  \node[label_node=cyan, above=0.25cm of layersXY1] {0.8};

  \node[gene_node , right=\separationofcircles cm of layersXY1] (layersYZ1) {Layers \\ yz plane};
  \node[label_node=magenta, above=0.25cm of layersYZ1] {0.4};

  \node[gene_node, right=\separationofcircles cm of layersYZ1] (layersXZ1) {Layers \\ xz plane};
  \node[label_node=violet!70, above=0.25cm of layersXZ1] {0.9};

  % Grouping gene nodes and numbers for container type 1
  \node[aux_gene_group, fit=(o1) (layersXZ1)] (group11) {};


% Place a dummy node above group1
\node[above=1cm of group11] (dummy) {};

% Extract the x and y coordinates from the node
\pgfpointanchor{dummy}{center} % You can replace 'center' with other anchor names (e.g., 'north', 'south west', etc.)
\pgfgetlastxy{\myNodeX}{\myNodeY}
% Now draw the orthotope using the dummy node's position
% Adjust the coordinates as needed

\node[above=of group11] (box11) {\tikz     \drawblock{1}{1}{1}{\lbo}{\wbo}{\hbo}{blue}{0}{0}{0};};
  \node[gene_group, fit=(box11) (group11)] (group1) {};
  % Container types
  \node[ellip, above=of group1] (containerType1) {Container type 1};

% After the definition of group1

% After the definition of group1

% Assuming the orthotope should be placed directly above the center of group1


% After the definition of group1


  % Gene nodes and numbers for container type 2
  \node[gene_node , right=1 cm of layersXZ1] (order2) {Order};
  \node[label_node=orange, above=0.25cm of order2] (o2) {0.2};

  \node[gene_node , right=\separationofcircles cm of order2] (orientation2) {Orientation};
  \node[label_node=lime, above=0.25cm of orientation2] {0.7};

  \node[gene_node , right=\separationofcircles cm of orientation2] (layersXY2) {Layers \\ xy plane};
  \node[label_node=cyan, above=0.25cm of layersXY2] {0.1};

  \node[gene_node , right=\separationofcircles cm of layersXY2] (layersYZ2) {Layers \\ yz plane};
  \node[label_node=magenta, above=0.25cm of layersYZ2] {0.1};

  \node[gene_node , right=\separationofcircles cm of layersYZ2] (layersXZ2) {Layers \\ xz plane};
  \node[label_node=violet!70, above=0.25cm of layersXZ2] {0.3};

  % Grouping gene nodes and numbers for container type 2
  \node[aux_gene_group, fit=(o2) (layersXZ2)] (group21) {};




% Place a dummy node above group1
\node[above=1cm of group21] (dummy1) {};

% Extract the x and y coordinates from the node
\pgfpointanchor{dummy1}{center} % You can replace 'center' with other anchor names (e.g., 'north', 'south west', etc.)
\pgfgetlastxy{\myNodeXX}{\myNodeYY}
% Now draw the orthotope using the dummy node's position
% Adjust the coordinates as needed
\node[above=of group21] (box21) {\tikz     \drawblock{1}{1}{1}{\lbt}{\wbt}{\hbt}{red}{0}{0}{0};};


  \node[gene_group, fit=(group21) (box21)] (group2) {};

  % First row: Set of Genes
  \node[box, left=of group1] (genes) {Genes};


\newcommand{\placeNodeAtIntersection}[3]{
  % #1 = new node name
  % #2 = node for x-coordinate
  % #3 = node for y-coordinate
  \path let \p1 = (#2), \p2 = (#3) in node (#1) at (\x1,\y2) {};
}

  \placeNodeAtIntersection{newNode}{genes}{containerType1};
  \node[box] at (newNode) {Set of Genes}; % This line actually creates the new node with the box style

\node[ellip, above=of group2] (containerType2) {Container type 2};

\node[below=2cm of group2] (fig2) {\tikz     \drawblock{1}{1}{1}{\lbo}{\wbo}{\hbo}{blue}{4}{2}{1};};
   
\node[below=2cm of group1] (fig1) {\tikz     \drawblock{1}{1}{1}{\wbt}{\hbt}{\lbt}{red}{2}{2}{1};};
\placeNodeAtIntersection{decodedgenes}{genes}{fig1};
\node[box] at (decodedgenes) {Decoded Genes}; % 
\coordinate (center13) at ($(group1.south)!1/3!(group1.south east)$);
\coordinate (center23) at ($(group2.south)!1/3!(group2.south west)$);
\path [line] (center13) -- (fig2.north west) ;
\path [line] (center23) -- (fig1.north east) ;


\end{tikzpicture}
